\subsection{Modular decomposition of pruning constraints}
\label{app:pruning-modules}

A nonempty block $B\subseteq W$ is a module of the allowed support
family $\mathcal F_r$ if two allowed supports that meet $B$ can
exchange their slices in $B$. That is, $(S\setminus B)\cup(T\cap B)$
is allowed whenever $S,T\in\mathcal F_r$ both meet $B$. A common
module satisfies the condition at every owner.

This definition gives a decomposition directly from the allowed
supports. Before computing the modules, delete any input absent
from all allowed supports of some owner, restrict the families to
sets of remaining inputs, and repeat. Every occurrence in a word
belongs to some pruning of each owner, so these deletions preserve
feasibility. If no input remains, feasibility fails. Otherwise
each family now supports every remaining input.

For an implication presentation, a candidate module can be checked
in polynomial time. Write $J_r(X)$ for the greatest allowed support
contained in $X\subseteq W$. To compute it, repeatedly delete a
head still present whose tail misses the current set. No deletion
can remove an input from an allowed support inside the current set,
so the final allowed support is the greatest one. A violation of
the module condition at the implication $j\Rightarrow\operatorname{hit}(T)$
requires the recipient to avoid $T\setminus B$ and the donor to
avoid $T\cap B$. Take their greatest possibilities
$J_r(W\setminus(T\setminus B))$ and $J_r(W\setminus(T\cap B))$.
They witness a violation exactly when both meet $B$ and the donor
contains $j$ if $j\in B$, or the recipient contains $j$ otherwise.
Testing each implication at each owner therefore suffices.

\begin{theorem}[Factorization over common modules]
\label{core:thm:pruning-module-factor}
Let $\mathcal P$ partition $W$ into common modules. For each block
$B$ use the inner families $\{S\cap B:S\in\mathcal F_r\}$. The
outer families on $\mathcal P$ are the images
$\{\{B:S\cap B\ne\varnothing\}:S\in\mathcal F_r\}$.
The original families are feasible exactly when the outer families
are feasible on the blocks with feasible inner families.
\end{theorem}
\begin{proof}
An allowed support is determined by an allowed active block set
and a choice of allowed nonempty slice at every active block.
For sufficiency of this condition, start with a support with the given
block set and replace each active slice using the module property.

Relabeling a word by blocks sends its pruning supports to their
images. If the word uses $B$, delete its leaves outside $B$ and
contract gates with just one remaining child. Every pruning support
of this restricted word is the $B$-slice of an original pruning
support. At a contracted gate owned by the pruning owner, reinstate
any pruning of the erased child, which has no leaves in $B$. At
a contracted foreign gate choose the retained child. Induction
gives the claim, so every block used by an outer witness is feasible.

Conversely substitute a feasible inner word at each leaf of an
outer witness. A pruning of the resulting word projects to the
pruning of the outer tree obtained by keeping its choices there.
Each active slice is the union of the inner pruning supports at
the selected occurrences of that block. It is allowed by union
closure, including when the block occurs more than once. The
support factorization then gives feasibility.
\end{proof}

Two subsets overlap if their intersection and both differences are
nonempty. A family of nonempty sets is \emph{partitive} if it contains
the whole set and every singleton, and is closed under intersection,
union, both differences and symmetric difference of overlapping
members.

\begin{theorem}[Common modules form a partitive family]
\label{core:thm:pruning-modules-partitive}
Suppose each $\mathcal F_r$ is closed under unions and every input
occurs in some member of each family. The common modules form a
partitive family.
\end{theorem}
\begin{proof}
Work at one owner. Write $A=B\setminus C$, $I=B\cap C$ and
$D=C\setminus B$ for overlapping modules. An $I$-slice swap between
two supports meeting $I$ can be done by first replacing the
$C$-slice, then taking the resulting $B$-slice back to the original
support. Thus $I$ is a module.

Choose a support meeting $I$. Any support meeting $B\cup C$ can
be made to meet $I$ by swapping its active $B$- or $C$-slice with
that of the chosen support. For a pair of supports now meeting $I$,
swapping first at $B$ and then at $C$ transports the entire
$(B\cup C)$-slice. Undo the donor's preparation by using its
original slice at the block used in the preparation. This proves
that $B\cup C$ is a module.

A support meeting $B$ can be made to meet $C$ while retaining its
$B$-slice. Add a support meeting $D$ by union, then restore the
$B$-slice. Given two supports $S,T$ meeting $B\cup C$, with $S$
meeting $A$, this prepares $S$ for a $C$-swap. Prepare $T$ in the
same way if it misses $C$. The $C$-swap gives a $B$-slice with
$S$'s trace on $A$ and $T$'s trace on $I$. If $T$ also meets $A$,
replace its $B$-slice by the one just obtained. This swaps $A$ alone.
Thus $A$ is a module, and dually so is $D$.

For the symmetric difference, suppose $S,T$ both meet $A\cup D$.
If $T$ meets $A$, achieve $T$'s traces on $A,D$ and $S$'s trace
on $I$ by constructing the $B$-slice with $T$'s $A$-trace and $S$'s
$I$-trace, and inserting it into $T$. If $T$ meets only $D$, apply
the dual construction. Now use the union module $B\cup C$ to
take the resulting slice into $S$. Singletons and the whole set
satisfy the module condition directly. The argument applies to
each owner, proving the common closure laws.
\end{proof}

A module is strong if it overlaps no module. The strong common
modules form a unique inclusion tree. At a node its maximal proper
strong submodules partition it. The decomposition theorem for
partitive families \cite{CheinHabibMaurer1981} says that the quotient
has either every nonempty set of child blocks as a module or just
the trivial modules. These are called complete and prime nodes.
Here the modules on the quotient are exactly the common modules
that are unions of child blocks, by the slice factorization in
\cref{core:thm:pruning-module-factor}. The module condition within a
node uses the families of slices there. It agrees with the original
condition on its subblocks because the entire node is a module.

\begin{proposition}[Complete nodes]
\label{core:prop:complete-pruning-modules}
On a fully supported domain, if every pair is a module of a
union-closed family $\mathcal F$, its nonempty members are exactly
the nonempty supersets of their common intersection $C$.
Consequently a complete node can be tested just on these compulsory
sets $C_r$ and the feasible children.
\end{proposition}
\begin{proof}
Given a nonempty member $S$ and a label $x$ outside it, choose
$y\in S$ and a donor containing $x$. Swap the pair $\{x,y\}$
into $S$, then union the result with $S$ to restore $y$ if needed.
This gives $S\cup\{x\}$, so the nonempty family is upward closed.
Now suppose $x\notin C$ and $y\ne x$ is another label in $S$.
Take a nonempty donor missing $x$ and add $y$ by upward closure.
Its singleton slice at $\{x,y\}$ lets us delete $x$ from $S$.
Thus every nonempty superset of $C$ may be obtained from the full
set by deleting the optional labels outside it.

At a complete node discard infeasible children. If none remain,
or a compulsory label has been discarded, the node fails. If at
most one owner has a nonempty compulsory set, a projection or that
owner's fold suffices. Otherwise, by the foreign singleton test of
\cref{core:prop:compulsory-prunings}, the nonempty compulsory sets
must all be the same singleton. Its projection then suffices.
\end{proof}

Thus modular decomposition reduces the general pruning problem to
tests on the feasible children of prime nodes. The following reduction
shows that prime quotients
retain the general recognition problem.

\begin{theorem}[Prime pruning instances]
\label{core:thm:prime-pruning-hardness}
On two or more owners, every finite pruning instance with at least
two inputs reduces to an equivalent instance with three times as
many inputs, fully supported at each owner, whose compulsory sets
are empty and whose only common modules are the trivial ones.
The reduction is polynomial for families presented by implications
$j\Rightarrow\operatorname{hit}(T)$. At every fixed $K\ge3$,
recognition for these prime instances is PSPACE-complete.
\end{theorem}
\begin{proof}
Write $x_1,\ldots,x_n$ for the original inputs and add
$D=\{a_i,b_i:1\le i\le n\}$. Choose two owners $r,s$. At $r$ pair
$a_i$ with $b_i$, and at $s$ pair $b_i$ with $a_{i+1}$ with indices
modulo $n$. For either owner a support meeting $D$ must take entire
pairs, and taking the $i$th pair requires $x_i$. At other owners
all supports meeting $D$ are permitted. Supports missing $D$
have their original conditions. These families are union closed.
Each input occurs in a member and is absent from another nonempty
member at every owner.

No feasible word uses $D$. To see this, retain only the two
pairing requirements and descend the word. At an $r$-gate,
both children inherit $s$'s pairing requirement. If an $r$-pruning
of one child took just one member of a pair, every $r$-pruning of
the other child would have to supply the missing label. Then a
singleton pruning at $s$ of the other child, as in
\cref{core:prop:compulsory-prunings}, violates its pairing
requirement. Hence both children also obey $r$'s requirement.
The $s$-gate is symmetric, and any other gate passes both
requirements to both children. A leaf in $D$ violates them.
Original words retain their legality, so feasibility is preserved
in both directions.

Suppose a common module $B$ contains more than one input and
contains an auxiliary label $a$. If one of its two partners $b$
lies outside $B$, use the support consisting of that pair and its
required original label. The corresponding owner also allows a
support containing another member of $B$ and missing the pair:
choose a different pair and include that member, using its own
pair if it is auxiliary. Swapping the slice into the first support
keeps $b$ and removes $a$. Thus $B$ contains both partners of
every auxiliary label it meets, and the two pairings together
force $D\subseteq B$. If $x_i$ lies outside $B$, swap the slice
of pair $i$ with its required $x_i$ into the support for a different
pair and its required label. The swap inserts pair $i$ but misses
$x_i$, a contradiction. Finally a block contained in the original
inputs and containing distinct $x_i,x_j$ fails to be a module:
swap the $x_j$-slice from the support for pair $j$ into the support
for pair $i$. This keeps pair $i$ and removes $x_i$.

For an implication presentation, adjoin $D$ to each old tail and
add the two implications enforcing equality of each pair and the
one requiring its original label. This gives the polynomial bound.
For hardness use the two positive singleton demands in
\cref{core:thm:two-color-pspace}, expressing each by monochromatic
omissions via \eqref{core:eq:total-monochromatic-source}. The
result is an unpointed pruning instance, to which the reduction
applies. For membership, try each possible input as an unposted
positive source and run the polynomial-space test on listed
omissions. Every word supplies such a source, so this suffices.
\end{proof}

For the graph families, a proper nonempty module in
a connected graph with an edge consists of independent vertices
with the same neighborhood. Indeed any crossing edge supplies a
singleton slice. Swapping it into an internal edge is impossible.
Swapping between crossing edges then forces all vertices of the
block to have the same neighbors. Conversely such a block admits
all slice swaps at the edge family. At the cover family a minimal
cover includes either all of it or none, since excluding a neighbor
forces every vertex of the block into the cover. So this
decomposition removes repeated neighborhoods, and in particular
leaves the seven-vertex example of
\cref{core:prop:nonrectangular-prunings} prime.
